% GENERATED FILE. Edit canonical lessons, then run npm run generate:books.
% canonical-source-hash: fnv1a64:cc3edcc69a41e9bd
% canonical-lessons: MR-C496-dir, MR-C496-nanand, MR-C496-mehuna, MR-C496-vahini, MR-C496-vyahi, MR-R496-review-490-496

\chapter{In-Laws}
\label{ch:mr-in-laws}
\hlchaptermodality{496}

\begin{chapteropening}
\textbf{By the end of this chapter:} \emph{I can name in-laws on a husband's or a wife's side.}

Complete the last lesson of chapter 496: I can name in-laws on a husband's or a wife's side.
\end{chapteropening}

\section[\texorpdfstring{\mr{दीर} (dīr)}{दीर (dīr)}]{\mr{दीर} (dīr) — a husband's brother}

\label{lesson:MR-C496-dir}



\begin{quote}
Before the new one: say the Marathi for a great-grandfather, then the Marathi for a father's sister.
\end{quote}

\subsection*{You'll want to know: \mr{दीर}}
\textbf{\mr{दीर}} — \emph{dīr} — \textquotedblleft{}a husband's brother\textquotedblright{}.

\begin{quote}
\mr{लग्नानंतर} \mr{तिचा} \mr{दीर} \mr{तिच्याशी} \mr{नीट} \mr{वागतो}.
\end{quote}

\begin{quote}
\emph{lagnānantar ticā dīr ticyāśī nīṭ vāgto.}
\end{quote}

\begin{quote}
After the wedding her husband's brother treats her well.
\end{quote}

Masculine.

A woman's husband's brother, usually the younger one.

\begin{cousinweb}
\textbf{\mr{दीर}} (\emph{dīr}) is a tadbhava, inherited from Sanskrit \emph{devṛ}, \textquotedblleft{}husband's brother\textquotedblright{}; CDIAL lists Marathi \emph{dīr} (CDIAL 6546).
\end{cousinweb}

\begin{grammarlens}[title={What you've built}]
One more word for describing what you lived and felt.
\end{grammarlens}

\subsection*{Your turn}
\begin{itemize}
\raggedright
  \item \emph{Say it:} \emph{dīr}
  \item \emph{Say it:} \emph{dīr}, once more
  \item \emph{Say it:} \emph{lagnānantar ticā dīr ticyāśī nīṭ vāgto.} — After the wedding her husband's brother treats her well.
  \item \emph{Recall:} say the Marathi for a great-grandfather, then the Marathi for a father's sister, then say \emph{dīr} again
\end{itemize}

\begin{tcolorbox}[breakable,colback=teal!4,colframe=teal!35!black,title={Before you move on}]
What does \mr{दीर} mean? (\textquotedblleft{}A husband's brother\textquotedblright{}.) Say it once more.
\end{tcolorbox}
\section[\texorpdfstring{\mr{नणंद} (naṇand)}{नणंद (naṇand)}]{\mr{नणंद} (naṇand) — a husband's sister}

\label{lesson:MR-C496-nanand}



\begin{quote}
Before the new one: say the Marathi for a father's sister, then the Marathi for a husband's brother.
\end{quote}

\subsection*{You'll want to know: \mr{नणंद}}
\textbf{\mr{नणंद}} — \emph{naṇand} — \textquotedblleft{}a husband's sister\textquotedblright{}.

\begin{quote}
\mr{सुनेची} \mr{आणि} \mr{नणंदेची} \mr{चांगली} \mr{मैत्री} \mr{आहे}.
\end{quote}

\begin{quote}
\emph{sunecī āṇi naṇandecī cāṅglī maitrī āhe.}
\end{quote}

\begin{quote}
The daughter-in-law and her husband's sister are good friends.
\end{quote}

Feminine.

Oblique \mr{नणंदे}-.

\begin{cousinweb}
\textbf{\mr{नणंद}} (\emph{naṇand}) is a tadbhava, inherited from Sanskrit \emph{nanāndṛ}, \textquotedblleft{}husband's sister\textquotedblright{}; CDIAL lists Marathi \emph{naṇãd} (CDIAL 6946).
\end{cousinweb}

\begin{grammarlens}[title={What you've built}]
One more word for describing what you lived and felt.
\end{grammarlens}

\subsection*{Your turn}
\begin{itemize}
\raggedright
  \item \emph{Say it:} \emph{naṇand}
  \item \emph{Say it:} \emph{naṇand}, once more
  \item \emph{Say it:} \emph{sunecī āṇi naṇandecī cāṅglī maitrī āhe.} — The daughter-in-law and her husband's sister are good friends.
  \item \emph{Recall:} say the Marathi for a father's sister, then the Marathi for a husband's brother, then say \emph{naṇand} again
\end{itemize}

\begin{tcolorbox}[breakable,colback=teal!4,colframe=teal!35!black,title={Before you move on}]
What does \mr{नणंद} mean? (\textquotedblleft{}A husband's sister\textquotedblright{}.) Say it once more.
\end{tcolorbox}
\section[\texorpdfstring{\mr{मेहुणा} (mehuṇā)}{मेहुणा (mehuṇā)}]{\mr{मेहुणा} (mehuṇā) — a wife's brother}

\label{lesson:MR-C496-mehuna}



\begin{quote}
Before the new one: say the Marathi for a husband's brother, then the Marathi for a husband's sister.
\end{quote}

\subsection*{You'll want to know: \mr{मेहुणा}}
\textbf{\mr{मेहुणा}} — \emph{mehuṇā} — \textquotedblleft{}a wife's brother\textquotedblright{}.

\begin{quote}
\mr{माझा} \mr{मेहुणा} \mr{मला} \mr{भेटायला} \mr{आला}.
\end{quote}

\begin{quote}
\emph{mājhā mehuṇā malā bheṭāylā ālā.}
\end{quote}

\begin{quote}
My wife's brother came to see me.
\end{quote}

Masculine.

Also a sister's husband; \mr{मेहुणी}, a wife's sister.

\begin{cousinweb}
\textbf{\mr{मेहुणा}} (\emph{mehuṇā}) is a tadbhava: CDIAL lists Marathi \emph{mehuṇā}, \textquotedblleft{}wife's brother, sister's husband\textquotedblright{}, under Sanskrit \emph{maithuna} (CDIAL 10341).
\end{cousinweb}

\begin{grammarlens}[title={What you've built}]
One more word for describing what you lived and felt.
\end{grammarlens}

\subsection*{Your turn}
\begin{itemize}
\raggedright
  \item \emph{Say it:} \emph{mehuṇā}
  \item \emph{Say it:} \emph{mehuṇā}, once more
  \item \emph{Say it:} \emph{mājhā mehuṇā malā bheṭāylā ālā.} — My wife's brother came to see me.
  \item \emph{Recall:} say the Marathi for a husband's brother, then the Marathi for a husband's sister, then say \emph{mehuṇā} again
\end{itemize}

\begin{tcolorbox}[breakable,colback=teal!4,colframe=teal!35!black,title={Before you move on}]
What does \mr{मेहुणा} mean? (\textquotedblleft{}A wife's brother\textquotedblright{}.) Say it once more.
\end{tcolorbox}
\section[\texorpdfstring{\mr{वहिनी} (vahinī)}{वहिनी (vahinī)}]{\mr{वहिनी} (vahinī) — a brother's wife}

\label{lesson:MR-C496-vahini}



\begin{quote}
Before the new one: say the Marathi for a husband's sister, then the Marathi for a wife's brother.
\end{quote}

\subsection*{You'll want to know: \mr{वहिनी}}
\textbf{\mr{वहिनी}} — \emph{vahinī} — \textquotedblleft{}a brother's wife\textquotedblright{}.

\begin{quote}
\mr{वहिनीने} \mr{सगळ्यांसाठी} \mr{स्वयंपाक} \mr{केला}.
\end{quote}

\begin{quote}
\emph{vahinīne sagḷyāṁsāṭhī svayampāk kelā.}
\end{quote}

\begin{quote}
My brother's wife cooked for everyone.
\end{quote}

Feminine.

Also a polite way to address a friend's wife.

\begin{cousinweb}
\textbf{\mr{वहिनी}} (\emph{vahinī}) is a tadbhava: CDIAL lists Marathi \emph{vahinī}, \textquotedblleft{}elder brother's wife\textquotedblright{}, under Sanskrit \emph{vadhū}, \textquotedblleft{}bride\textquotedblright{} (CDIAL 11250).
\end{cousinweb}

\begin{grammarlens}[title={What you've built}]
One more word for describing what you lived and felt.
\end{grammarlens}

\subsection*{Your turn}
\begin{itemize}
\raggedright
  \item \emph{Say it:} \emph{vahinī}
  \item \emph{Say it:} \emph{vahinī}, once more
  \item \emph{Say it:} \emph{vahinīne sagḷyāṁsāṭhī svayampāk kelā.} — My brother's wife cooked for everyone.
  \item \emph{Recall:} say the Marathi for a husband's sister, then the Marathi for a wife's brother, then say \emph{vahinī} again
\end{itemize}

\begin{tcolorbox}[breakable,colback=teal!4,colframe=teal!35!black,title={Before you move on}]
What does \mr{वहिनी} mean? (\textquotedblleft{}A brother's wife\textquotedblright{}.) Say it once more.
\end{tcolorbox}
\section[\texorpdfstring{\mr{व्याही} (vyāhī)}{व्याही (vyāhī)}]{\mr{व्याही} (vyāhī) — a child's father-in-law}

\label{lesson:MR-C496-vyahi}



\begin{quote}
Before the new one: say the Marathi for a wife's brother, then the Marathi for a brother's wife.
\end{quote}

\subsection*{You'll want to know: \mr{व्याही}}
\textbf{\mr{व्याही}} — \emph{vyāhī} — \textquotedblleft{}a child's father-in-law\textquotedblright{}.

\begin{quote}
\mr{लग्नात} \mr{दोन्ही} \mr{व्याही} \mr{आनंदाने} \mr{भेटले}.
\end{quote}

\begin{quote}
\emph{lagnāt donhī vyāhī ānandāne bheṭle.}
\end{quote}

\begin{quote}
At the wedding the two fathers-in-law met happily.
\end{quote}

Masculine.

\mr{विहीण}, a child's mother-in-law: the parents of a married couple are each other's \mr{व्याही}.

\begin{cousinweb}
\textbf{\mr{व्याही}} (\emph{vyāhī}) is a tadbhava: CDIAL lists Marathi \emph{vyāhī}, \textquotedblleft{}son's or daughter's father-in-law\textquotedblright{}, under Sanskrit \emph{vivāhin} (CDIAL 11925), a word of marriage.
\end{cousinweb}

\begin{grammarlens}[title={What you've built}]
One more word for describing what you lived and felt.
\end{grammarlens}

\subsection*{Your turn}
\begin{itemize}
\raggedright
  \item \emph{Say it:} \emph{vyāhī}
  \item \emph{Say it:} \emph{vyāhī}, once more
  \item \emph{Say it:} \emph{lagnāt donhī vyāhī ānandāne bheṭle.} — At the wedding the two fathers-in-law met happily.
  \item \emph{Recall:} say the Marathi for a wife's brother, then the Marathi for a brother's wife, then say \emph{vyāhī} again
\end{itemize}

\begin{tcolorbox}[breakable,colback=teal!4,colframe=teal!35!black,title={Before you move on}]
What does \mr{व्याही} mean? (\textquotedblleft{}A child's father-in-law\textquotedblright{}.) Say it once more.
\end{tcolorbox}
\section[(dialogue)]{Looking back: from Staying Healthy to In-Laws}

\label{lesson:MR-R496-review-490-496}



\begin{quote}
Say the Marathi for a diet regimen and healthy.
\end{quote}

\subsection*{Your turn}
Say these aloud:

\begin{itemize}
\raggedright
  \item a diet regimen, healthy, immunity, to apply heat, to rub
  \item insomnia, to snore, to yawn, to take a nap, to lie down
  \item to have hiccups, to burp, to spit, to get breathless, to choke on food
  \item a crutch, to feel dizzy, to stumble, to shiver with cold, to limp
  \item a relationship, a relative, a household, an orphan, full-blood
  \item a brother's son, a sister's son, a grandson, a great-grandfather, a father's sister
  \item a husband's brother, a husband's sister, a wife's brother, a brother's wife, a child's father-in-law
\end{itemize}

\begin{tcolorbox}[breakable,colback=teal!4,colframe=teal!35!black,title={Before you move on}]
A diet regimen? (\textbf{\mr{पथ्य}}, \emph{pathya}.) Healthy? (\textbf{\mr{निरोगी}}, \emph{nirogī}.) Immunity? (\textbf{\mr{प्रतिकारशक्ती}}, \emph{pratikārśaktī}.) To apply heat? (\textbf{\mr{शेकणे}}, \emph{śekṇe}.) To rub? (\textbf{\mr{चोळणे}}, \emph{coḷṇe}.) Insomnia? (\textbf{\mr{निद्रानाश}}, \emph{nidrānāś}.) To snore? (\textbf{\mr{घोरणे}}, \emph{ghorṇe}.) To yawn? (\textbf{\mr{जांभई} \mr{देणे}}, \emph{jāmbhaī deṇe}.) To take a nap? (\textbf{\mr{डुलकी} \mr{घेणे}}, \emph{ḍulkī gheṇe}.) To lie down? (\textbf{\mr{निजणे}}, \emph{nijṇe}.) To have hiccups? (\textbf{\mr{उचकी} \mr{लागणे}}, \emph{uckī lāgṇe}.) To burp? (\textbf{\mr{ढेकर} \mr{देणे}}, \emph{ḍhekar deṇe}.) To spit? (\textbf{\mr{थुंकणे}}, \emph{thuṅkṇe}.) To get breathless? (\textbf{\mr{धाप} \mr{लागणे}}, \emph{dhāp lāgṇe}.) To choke on food? (\textbf{\mr{ठसका} \mr{लागणे}}, \emph{ṭhaskā lāgṇe}.) A crutch? (\textbf{\mr{कुबडी}}, \emph{kubḍī}.) To feel dizzy? (\textbf{\mr{चक्कर} \mr{येणे}}, \emph{cakkar yeṇe}.) To stumble? (\textbf{\mr{अडखळणे}}, \emph{aḍkhaḷṇe}.) To shiver with cold? (\textbf{\mr{कुडकुडणे}}, \emph{kuḍkuḍṇe}.) To limp? (\textbf{\mr{लंगडणे}}, \emph{laṅgaḍṇe}.) A relationship? (\textbf{\mr{नाते}}, \emph{nāte}.) A relative? (\textbf{\mr{नातेवाईक}}, \emph{nātevāīk}.) A household? (\textbf{\mr{परिवार}}, \emph{parivār}.) An orphan? (\textbf{\mr{अनाथ}}, \emph{anāth}.) Full-blood? (\textbf{\mr{सख्खा}}, \emph{sakkhā}.) A brother's son? (\textbf{\mr{पुतण्या}}, \emph{putaṇyā}.) A sister's son? (\textbf{\mr{भाचा}}, \emph{bhācā}.) A grandson? (\textbf{\mr{नातू}}, \emph{nātū}.) A great-grandfather? (\textbf{\mr{पणजोबा}}, \emph{paṇjobā}.) A father's sister? (\textbf{\mr{आत्या}}, \emph{ātyā}.) A husband's brother? (\textbf{\mr{दीर}}, \emph{dīr}.) A husband's sister? (\textbf{\mr{नणंद}}, \emph{naṇand}.) A wife's brother? (\textbf{\mr{मेहुणा}}, \emph{mehuṇā}.) A brother's wife? (\textbf{\mr{वहिनी}}, \emph{vahinī}.) A child's father-in-law? (\textbf{\mr{व्याही}}, \emph{vyāhī}.)
\end{tcolorbox}
